\chapter{Linear Combinations, Span, and Geometry}

In the last chapter you met a vector as three things at once: a list of numbers, a point in space, and an arrow from the origin. Now we put scaling and addition together into the central operation of linear algebra: the \emph{linear combination}.

Much of this subject comes back to one question: \emph{starting from a few vectors, where can I get to by scaling them and adding them up?} The set of reachable destinations is called the \textbf{span}. It is always a point, a line, a plane, all of space, or a higher-dimensional version of these, and these shapes are what lies behind every system of equations and every matrix.

\chapterline

\section{Scaling and Adding: The Two Moves}

This chapter uses only two operations on vectors, and you met both in the last chapter.

\textbf{Move one: scaling.} Take a vector $\vec{v}$ and a number $c$ (a \emph{scalar}). Then $c\,\vec{v}$ stretches $\vec{v}$ by the factor $c$. If $c>1$ it grows; if $0<c<1$ it shrinks; if $c<0$ it flips to point the other way; if $c=0$ it collapses to the origin. Scaling never lets $\vec{v}$ leave its own line: whatever $c$ you pick, $c\,\vec{v}$ sits somewhere on the infinite line through the origin and the tip of $\vec{v}$ (Figure~\ref{ch03:scaling-line}).

\textbf{Move two: adding.} Take two vectors and add them tip-to-tail (the parallelogram rule from the last chapter). Addition lets you leave a single line and explore sideways.

\begin{intuitionbox}
Scaling moves you \emph{along} a direction. Adding lets you \emph{mix} directions. A linear combination does both: pick how much of each direction you want (scaling), then mix them (adding). The rest of the chapter watches what shapes you can reach by doing this.
\end{intuitionbox}

\begin{figure}[H]
\centering
\begin{tikzpicture}[scale=1.0,>=Stealth]
  \draw[help lines,gray!30] (-3.4,-1.4) grid (3.4,3.4);
  \draw[->] (-3.4,0)--(3.4,0) node[right]{$x$};
  \draw[->] (0,-1.4)--(0,3.4) node[above]{$y$};
  % base vector v = (1,1)
  \draw[->,vecblue,very thick] (0,0)--(1,1) node[midway,above left]{$\vec{v}$};
  % scaled copies
  \draw[->,vecblue,thick,dashed] (0,0)--(2,2) node[above left]{$2\vec{v}$};
  \draw[->,vecblue,thick,dashed] (0,0)--(-1,-1) node[above left]{$-\vec{v}$};
  % the whole line
  \draw[vecblue!50,thick] (-1.3,-1.3)--(3.3,3.3);
  \node[vecblue!70!black,anchor=north west] at (2.4,2.0) {\small the line $c\vec{v}$};
\end{tikzpicture}
\caption{Scaling alone. Every scalar multiple $c\,\vec{v}$ of a single nonzero vector lands on one line through the origin. Scaling can stretch, shrink, or flip, but it cannot escape the line.}
\label{ch03:scaling-line}
\end{figure}

\section{The Linear Combination}

Combine the two moves into one definition. Take vectors $\vec{v}_1, \vec{v}_2, \dots, \vec{v}_k$ and scalars $c_1, c_2, \dots, c_k$. The expression
\[
  c_1 \vec{v}_1 + c_2 \vec{v}_2 + \cdots + c_k \vec{v}_k
\]
is called a \textbf{linear combination} of the vectors. The scalars $c_i$ are the \textbf{weights} (or coefficients). The word ``linear'' means that the only operations allowed are scaling and adding: no squaring a vector, no multiplying two vectors together, no taking a sine.

In the definition below, $k$ is the number of vectors we combine, and $\vec{v}_1,\dots,\vec{v}_k \in \R^n$ says that each of them is a list of $n$ real numbers.

\begin{definition}[Linear combination]\label{ch03:def-lincomb}
A vector $\vec{b}$ is a \textbf{linear combination} of $\vec{v}_1,\dots,\vec{v}_k \in \R^n$ if there exist scalars $c_1,\dots,c_k \in \R$ such that
\[
  \vec{b} = c_1\vec{v}_1 + c_2\vec{v}_2 + \cdots + c_k\vec{v}_k.
\]
\end{definition}

Make it concrete. Take two vectors in the plane,
\[
  \vec{v}_1 = \colvec{2\\1}, \qquad \vec{v}_2 = \colvec{1\\3}.
\]
Here is the linear combination with weights $c_1 = 2$ and $c_2 = 1$:
\[
  2\colvec{2\\1} + 1\colvec{1\\3}
  = \colvec{4\\2} + \colvec{1\\3}
  = \colvec{5\\5}.
\]
And with weights $c_1 = 1$, $c_2 = -1$:
\[
  1\colvec{2\\1} - 1\colvec{1\\3}
  = \colvec{2\\1} - \colvec{1\\3}
  = \colvec{1\\-2}.
\]
Each choice of weights gives a destination, and different weights give different destinations. The central question of the chapter: as we sweep the weights over \emph{all} possible real numbers, what is the full set of destinations?

\begin{figure}[H]
\centering
\begin{tikzpicture}[scale=0.75,>=Stealth]
  \draw[help lines,gray!25] (-1.4,-2.4) grid (5.9,5.9);
  \draw[->] (-1.4,0)--(5.9,0) node[right]{$x$};
  \draw[->] (0,-2.4)--(0,5.9) node[above]{$y$};
  % v1, v2
  \draw[->,vecblue,very thick] (0,0)--(2,1) node[below right]{$\vec{v}_1$};
  \draw[->,vecred,very thick] (0,0)--(1,3) node[above left]{$\vec{v}_2$};
  % build 2v1 + v2
  \draw[->,vecblue,thick,dashed] (0,0)--(4,2) node[below right]{$2\vec{v}_1$};
  \draw[->,vecred,thick,dashed] (4,2)--(5,5) node[midway,below right]{$\vec{v}_2$};
  \draw[->,targetgreen,very thick] (0,0)--(5,5) node[above right]{$2\vec{v}_1+\vec{v}_2$};
  \fill[targetgreen] (5,5) circle (2.5pt);
\end{tikzpicture}
\caption{The linear combination $2\vec{v}_1 + \vec{v}_2 = \colvec{5\\5}$ built tip-to-tail: walk out $2\vec{v}_1$ (blue, doubled), then add $\vec{v}_2$ (red) from there. The green arrow is the result.}
\label{ch03:one-combo}
\end{figure}

\begin{focusbox}
A linear combination is a weighted sum: each scalar determines how much of a direction to
include, and their vector sum is the result. The span (defined next) is the collection of all
points reachable by varying these weights across all real values.
\end{focusbox}

\textbf{Quick mental check.} Suppose your building blocks are the standard coordinate vectors $\vec{e}_1 = \colvec{1\\0}$ and $\vec{e}_2 = \colvec{0\\1}$. What weights do you need to reach $\colvec{7\\-4}$? You do not have to solve any equations: $7\vec{e}_1 - 4\vec{e}_2 = \colvec{7\\-4}$. The coordinates of a vector \emph{are} its weights in the recipe of the standard axes.

\section{Span: The Set of All Combinations}

\begin{definition}[Span]\label{ch03:def-span}
The \textbf{span} of vectors $\vec{v}_1,\dots,\vec{v}_k$, written $\spn\{\vec{v}_1,\dots,\vec{v}_k\}$, is the set of \emph{all} their linear combinations:
\[
  \spn\{\vec{v}_1,\dots,\vec{v}_k\}
  = \set{\, c_1\vec{v}_1 + \cdots + c_k\vec{v}_k : c_1,\dots,c_k \in \R \,}.
\]
\end{definition}

The braces read ``the set of all vectors of the form $c_1\vec{v}_1 + \cdots + c_k\vec{v}_k$, where (the colon) the weights $c_1,\dots,c_k$ range over all real numbers.''

Read the word ``span'' as a verb and a noun. As a verb: the vectors \emph{span} a region -- they reach it, they cover it. As a noun: the span \emph{is} that region, the complete set of reachable points. The weights are the dials; the span is everywhere the dials can take you.

A span is never a random scatter of points. It is always one of a short list of clean geometric shapes: the origin alone, a line, a plane, all of $3$-space, and so on up. We will walk up that ladder one rung at a time.

\subsection{Span of one vector: a line (usually)}

Start with a single nonzero vector $\vec{v}$. Its span is the set of all $c\,\vec{v}$. We already saw this in Figure~\ref{ch03:scaling-line}: scaling one vector traces out a line through the origin. So
\[
  \spn\{\vec{v}\} = \text{the line through the origin in the direction of } \vec{v}.
\]

There is one exception, and it matters. If $\vec{v} = \vec{0}$, then $c\,\vec{0} = \vec{0}$ for every $c$. The span of the zero vector is the origin alone, the one-element set $\{\vec{0}\}$.

\begin{trapbox}
``A line'' here always means a line \emph{through the origin}. The span of one vector can never be a line that misses the origin, because $c=0$ always produces $\vec{0}$, so the origin is in \emph{every} span. If someone hands you a line that does not pass through the origin and asks ``is this a span?'' the answer is no. (It is a \emph{shifted} line, a line through the origin slid sideways; such lines are called \emph{affine} lines, and we will meet them when we solve $\mat{A}\vec{x}=\vec{b}$ for nonzero $\vec{b}$.)
\end{trapbox}

\textbf{Quick mental check.} Is the origin $\vec{0}$ in the span of \emph{any} collection of vectors $\vec{v}_1,\dots,\vec{v}_k$? Yes, always: set every weight to zero ($0\vec{v}_1 + \cdots + 0\vec{v}_k = \vec{0}$). No matter what vectors you choose, their span is anchored at the origin.

\subsection{Span of two vectors: a plane (usually)}

Now take two vectors $\vec{v}_1$ and $\vec{v}_2$ and ask for $\spn\{\vec{v}_1,\vec{v}_2\} = \set{c_1\vec{v}_1 + c_2\vec{v}_2}$.

Think of it as a grid. Sliding $c_1$ moves you along the $\vec{v}_1$ direction; sliding $c_2$ moves you along the $\vec{v}_2$ direction. With two separate dials you can reach any point on a flat sheet -- a \textbf{plane} through the origin. The two vectors lay down a (possibly skewed) coordinate grid, and the span is the whole infinite sheet that grid lives in.

\begin{figure}[H]
\centering
\begin{tikzpicture}[scale=0.85,>=Stealth]
  \draw[help lines,gray!20] (-2.4,-2.4) grid (4.4,4.4);
  \draw[->] (-2.4,0)--(4.4,0) node[right]{$x$};
  \draw[->] (0,-2.4)--(0,4.4) node[above]{$y$};
  % skewed grid from v1=(2,1), v2=(1,2)
  \foreach \c in {-2,-1,0,1,2}{
    \draw[vecblue!30] ($\c*(2,1)+(-2,-4)$)--($\c*(2,1)+(2,4)$);
    \draw[vecred!30]  ($\c*(1,2)+(-4,-2)$)--($\c*(1,2)+(4,2)$);
  }
  \draw[->,vecblue,very thick] (0,0)--(2,1) node[below right]{$\vec{v}_1$};
  \draw[->,vecred,very thick]  (0,0)--(1,2) node[above left]{$\vec{v}_2$};
  \fill[targetgreen] (3,3) circle (2.5pt) node[above right,targetgreen]{\small $\vec{v}_1+\vec{v}_2$};
\end{tikzpicture}
\caption{Two non-parallel vectors lay down a skewed grid. Sliding the two weights $c_1,c_2$ reaches every point of the plane -- here, in $\R^2$, that plane is the entire flat page. The blue and red families of lines are the ``coordinate lines'' of the grid.}
\label{ch03:two-grid}
\end{figure}

In $\R^2$ that plane \emph{is} all of $\R^2$: two non-parallel vectors in the plane span the whole plane (Theorem~\ref{ch03:thm-span-r2} below proves this). But the same two vectors living in $\R^3$ span a genuine $2$-dimensional plane floating inside $3$-space -- a flat sheet through the origin, with all of $3$-space stretching above and below it that you cannot reach.

\textbf{The exception again.} What if $\vec{v}_2$ happens to lie on the same line as $\vec{v}_1$ -- say $\vec{v}_2 = 3\vec{v}_1$? Then the second dial gives you nothing new. Anything you build, $c_1\vec{v}_1 + c_2(3\vec{v}_1) = (c_1 + 3c_2)\vec{v}_1$, is still a multiple of $\vec{v}_1$. The span collapses back down to a single line. Two vectors span a plane only when they point in different directions.

\subsection{Span of three vectors: all of space (usually)}

Three vectors in $\R^3$ that point in three genuinely different directions (none lies in the plane of the other two) span all of $\R^3$: with three dials you can reach every point in space.

The exceptions now stack up. They are what the word \emph{dependence} will describe:

\begin{itemize}
  \item If the third vector $\vec{v}_3$ already lies in the plane spanned by $\vec{v}_1$ and $\vec{v}_2$, it adds no new direction. The span stays a plane.
  \item If all three lie on one line, the span is just that line.
  \item If all three are $\vec{0}$, the span is the origin.
\end{itemize}

\begin{focusbox}
The ladder of spans: \emph{point, line, plane, space, \dots} Each new vector either bumps you up one rung (it points somewhere new) or it does nothing (it was already reachable). That choice between ``new direction'' and ``nothing new'' is the seed of \emph{linear independence}, named later in this chapter and studied fully in a coming chapter. Figure~\ref{ch03:ladder} draws the ladder.
\end{focusbox}

\begin{figure}[htbp]
\centering
\begin{tikzpicture}[>=Stealth,
   rung/.style={draw=spacegray,rounded corners=3pt,minimum width=2.6cm,minimum height=2.1cm,fill=gray!4}]
  % four rungs
  \node[rung] (P) at (0,0) {};
  \node[rung] (L) at (3.4,0) {};
  \node[rung] (S) at (6.8,0) {};
  \node[rung] (V) at (10.2,0) {};
  % point
  \fill[vecblue] (0,0.1) circle (2.5pt);
  \node[below] at (0,-0.45) {\small point $\{\vec{0}\}$};
  % line
  \draw[vecblue!60,thick] (2.6,-0.35)--(4.2,0.55);
  \fill[black] (3.4,0.1) circle (1.2pt);
  \node[below] at (3.4,-0.45) {\small line};
  % plane
  \fill[targetgreen!20] (6.0,-0.3)--(7.3,-0.3)--(7.6,0.5)--(6.3,0.5)--cycle;
  \draw[targetgreen!60!black] (6.0,-0.3)--(7.3,-0.3)--(7.6,0.5)--(6.3,0.5)--cycle;
  \fill[black] (6.8,0.1) circle (1.2pt);
  \node[below] at (6.8,-0.45) {\small plane};
  % space
  \draw[gray] (9.6,-0.3) rectangle (10.6,0.5);
  \draw[gray] (9.6,0.5)--(9.85,0.7)--(10.85,0.7)--(10.6,0.5);
  \draw[gray] (10.85,0.7)--(10.85,-0.1)--(10.6,-0.3);
  \fill[structurepurple!15] (9.6,-0.3) rectangle (10.6,0.5);
  \node[below] at (10.2,-0.45) {\small all of $\R^3$};
  % dimension labels
  \foreach \x/\d in {0/0,3.4/1,6.8/2,10.2/3}{
    \node[spacegray] at (\x,1.3) {\small dimension $\d$};
  }
  % arrows: new direction
  \foreach \a/\b in {P/L,L/S,S/V}{
    \draw[->,thick,targetgreen!60!black] (\a.north east)++(0,-0.35) -- ($(\b.north west)+(0,-0.35)$);
  }
  \node[targetgreen!50!black,align=center] at (5.1,-1.75) {\small a vector pointing somewhere new: up one rung (green arrows)};
  % loops: nothing new
  \foreach \a in {P,L,S,V}{
    \draw[->,thick,trapred] ($(\a.south)+(0.3,-0.05)$) arc (0:-180:0.3);
  }
  \node[trapred] at (5.1,-2.25) {\small a vector already in the span: stay on the same rung (red loops)};
\end{tikzpicture}
\caption{The ladder of spans in $\R^3$. Each vector you add either points in a new direction and lifts the span one rung (green arrows), or already lies in the span and leaves it unchanged (red loops). No single vector can lift the span two rungs at once.}
\label{ch03:ladder}
\end{figure}

\begin{figure}[H]
\centering
\begin{tikzpicture}[scale=1.5,>=Stealth,
   x={(0.9cm,0.32cm)}, y={(-0.86cm,0.42cm)}, z={(0cm,1cm)}]
  % a plane through origin spanned by v1, v2, with a third vector poking out
  \fill[targetgreen!12] (-1.2,-1.2,0)--(1.4,-1.2,0)--(1.4,1.2,0)--(-1.2,1.2,0)--cycle;
  \draw[gray!50] (-1.2,-1.2,0)--(1.4,-1.2,0)--(1.4,1.2,0)--(-1.2,1.2,0)--cycle;
  % axes hint
  \draw[->,gray] (0,0,0)--(1.6,0,0) node[right]{\small $x$};
  \draw[->,gray] (0,0,0)--(0,1.5,0) node[left]{\small $y$};
  \draw[->,gray] (0,0,0)--(0,0,1.5) node[above]{\small $z$};
  % v1, v2 in plane
  \draw[->,vecblue,very thick] (0,0,0)--(1,0.3,0) node[right]{$\vec{v}_1$};
  \draw[->,vecred,very thick]  (0,0,0)--(0.2,1,0) node[below left]{$\vec{v}_2$};
  % v3 out of plane
  \draw[->,structurepurple,very thick] (0,0,0)--(0.3,0.2,1) node[above right]{$\vec{v}_3$};
  \node[targetgreen!60!black] at (-0.7,0.9,0) {\small $\spn\{\vec{v}_1,\vec{v}_2\}$};
\end{tikzpicture}
\caption{In $\R^3$: $\vec{v}_1,\vec{v}_2$ span the shaded plane through the origin. A third vector $\vec{v}_3$ that pokes \emph{out} of that plane is a brand new direction -- together the three span all of $\R^3$. If $\vec{v}_3$ had lain flat in the plane, it would have added nothing.}
\label{ch03:three-3d}
\end{figure}

\begin{walkbox}
Picture a single vector in $\R^3$, drawn as an arrow. Its span is a line. Now picture sweeping that line around -- how many \emph{different} lines through the origin are there? Infinitely many, one for each direction. Now: how many different \emph{planes} through the origin are there in $\R^3$? Also infinitely many. A span in $\R^3$ is therefore determined by two things: \emph{its dimension and its orientation}. On your walk, try to feel the difference between ``what dimension is this span'' and ``which way is it tilted.''
\end{walkbox}

\section{A Fully Worked Span Question}

Here is one span question worked by hand from start to finish.

\begin{example}[Is $\vec{b}$ in the span?]\label{ch03:ex-in-span}
Let
\[
  \vec{v}_1 = \colvec{1\\2}, \qquad \vec{v}_2 = \colvec{3\\1}, \qquad \vec{b} = \colvec{5\\0}.
\]
Is $\vec{b} \in \spn\{\vec{v}_1, \vec{v}_2\}$? In words: can we hit the target $\vec{b}$ by mixing $\vec{v}_1$ and $\vec{v}_2$?
\end{example}

Being in the span means there \emph{exist} weights $c_1, c_2$ with
\[
  c_1\colvec{1\\2} + c_2\colvec{3\\1} = \colvec{5\\0}.
\]
Reading off each coordinate turns this single vector equation into two scalar equations -- one per row:
\[
  \begin{aligned}
    c_1 + 3c_2 &= 5 \quad (\text{top row}),\\
    2c_1 + c_2 &= 0 \quad (\text{bottom row}).
  \end{aligned}
\]
From the second equation, $c_2 = -2c_1$. Substitute into the first:
\begin{align*}
  c_1 + 3(-2c_1) &= 5\\
  -5c_1 &= 5\\
  c_1 &= -1.
\end{align*}
Then $c_2 = -2(-1) = 2$. Now \emph{check} -- always check:
\[
  -1\colvec{1\\2} + 2\colvec{3\\1}
  = \colvec{-1\\-2} + \colvec{6\\2}
  = \colvec{5\\0} = \vec{b}. \checkmark
\]
Yes. $\vec{b}$ is in the span, reached with weights $(-1, 2)$. And because $\vec{v}_1, \vec{v}_2$ point in different directions, they span \emph{all} of $\R^2$ (Theorem~\ref{ch03:thm-span-r2}), so this would have worked for any target $\vec{b}$ we picked.

\begin{focusbox}
``Is $\vec{b}$ in the span?'' is the same question as ``does this system of linear equations have a solution?'' That link connects the geometry of spans and planes to the algebra of linear systems, and we use it throughout the next several chapters.
\end{focusbox}

\section{The Column Picture: $\mat{A}\vec{x}$ as a Combination of Columns}

Strang opens his course with the idea in this section: a matrix times a vector \emph{is} a linear combination. Chapter~1 stated it; here we derive it.

Take a matrix and a vector:
\[
  \mat{A} = \begin{bmatrix} 1 & 3 \\ 2 & 1 \end{bmatrix}, \qquad \vec{x} = \colvec{c_1 \\ c_2}.
\]
The usual ``rows times the vector'' recipe for $\mat{A}\vec{x}$ goes row by row: multiply each entry of a row by the matching entry of $\vec{x}$ and add the products; each row gives one entry of the answer. Here that gives
\[
  \mat{A}\vec{x}
  = \begin{bmatrix} 1 & 3 \\ 2 & 1 \end{bmatrix}\colvec{c_1\\c_2}
  = \colvec{1\cdot c_1 + 3\cdot c_2 \\ 2\cdot c_1 + 1\cdot c_2}.
\]
Now \emph{regroup that right-hand side by columns} instead of by rows:
\[
  \colvec{c_1 + 3c_2 \\ 2c_1 + c_2}
  = \colvec{c_1 \\ 2c_1} + \colvec{3c_2 \\ c_2}
  = c_1\colvec{1\\2} + c_2\colvec{3\\1}.
\]
Those two final vectors are the columns of $\mat{A}$: $\colvec{1\\2}$ is the first column, $\colvec{3\\1}$ the second. So
\[
  \mat{A}\vec{x} = c_1\,(\text{column 1 of }\mat{A}) + c_2\,(\text{column 2 of }\mat{A}).
\]

\begin{formalbox}
\textbf{The column rule for $\mat{A}\vec{x}$.} Let $\mat{A}$ be an $m\times n$ matrix ($m$ rows, $n$ columns). If $\mat{A}$ has columns $\vec{a}_1,\vec{a}_2,\dots,\vec{a}_n$ (each a vector in $\R^m$) and $\vec{x} = \colvec{x_1\\ \vdots \\ x_n}$, then
\[
  \mat{A}\vec{x} = x_1\vec{a}_1 + x_2\vec{a}_2 + \cdots + x_n\vec{a}_n.
\]
\textbf{$\mat{A}\vec{x}$ is the linear combination of the columns of $\mat{A}$ with weights taken from $\vec{x}$.}
\end{formalbox}

This is the \textbf{column picture}. The entries of $\vec{x}$ are the \emph{weights} in a linear combination, and the things being combined are the columns.

Figure~\ref{ch03:column-picture} draws this for the system we solved by hand in Example~\ref{ch03:ex-in-span}: the columns of $\mat{A}$ are the two arrows you are allowed to scale, and $\mat{A}\vec{x}=\vec{b}$ asks you to dial the weights until tip-to-tail you land exactly on the target $\vec{b}$.

\begin{figure}[H]
\centering
\begin{tikzpicture}[scale=0.7,>=Stealth]
  \draw[help lines,gray!22] (-2.4,-2.9) grid (6.4,2.9);
  \draw[->] (-2.4,0)--(6.4,0) node[right]{$x$};
  \draw[->] (0,-2.9)--(0,2.9) node[above]{$y$};
  % columns a1=(1,2), a2=(3,1)
  \draw[->,vecblue,very thick] (0,0)--(1,2) node[above left]{$\vec{a}_1$};
  \draw[->,vecred,very thick] (0,0)--(3,1) node[above=2pt]{$\vec{a}_2$};
  % combination -1 a1 + 2 a2 = (5,0): walk -a1 then +2a2
  \draw[->,vecblue,thick,dashed] (0,0)--(-1,-2) node[below left]{$-\vec{a}_1$};
  \draw[->,vecred,thick,dashed] (-1,-2)--(5,0);
  \node[vecred!70!black] at (2.3,-1.5) {\small $2\vec{a}_2$};
  \draw[->,targetgreen,very thick] (0,0)--(5,0);
  \fill[targetgreen] (5,0) circle (2.8pt) node[above=4pt,targetgreen!60!black]{\small $\vec{b}=\colvec{5\\0}$};
\end{tikzpicture}
\caption{The column picture of $\mat{A}\vec{x}=\vec{b}$ for $\mat{A}=\protect\begin{bmatrix}1&3\\2&1\protect\end{bmatrix}$, $\vec{b}=\colvec{5\\0}$. Solving the system means choosing weights so that $x_1\vec{a}_1+x_2\vec{a}_2$ reaches the green target. The weights $(-1,2)$ from Example~\ref{ch03:ex-in-span} do it: walk $-\vec{a}_1$ (down-left), then add $2\vec{a}_2$, and you land on $\vec{b}$.}
\label{ch03:column-picture}
\end{figure}

\begin{intuitionbox}
There are two ways to read $\mat{A}\vec{x}$, and it pays to keep both in mind:
\begin{itemize}
  \item \textbf{Row picture:} each entry of the output is computed row by row as a dot product (multiplying matching entries across the row and down $\vec{x}$, then adding them up). Good for computing one number at a time, and for thinking of each equation as a line or plane.
  \item \textbf{Column picture:} the whole output is a linear combination of the columns, with weights $\vec{x}$. Good for thinking about \emph{which outputs are reachable at all}.
\end{itemize}
The row picture answers ``what is the answer?'' The column picture answers ``which answers are even possible?'' This chapter lives in the column picture.
\end{intuitionbox}

\begin{labbox}
Do this one with pen and paper before reading on. Let
\[
  \mat{A} = \begin{bmatrix} 2 & 0 & 1 \\ 1 & 3 & 1 \end{bmatrix}, \qquad \vec{x} = \colvec{1\\2\\-1}.
\]
Compute $\mat{A}\vec{x}$ \emph{twice}: once the row way (multiplying matching entries across each row and adding), once the column way ($1$ times column one, plus $2$ times column two, minus $1$ times column three, added up). You should get the same vector, $\colvec{1\\6}$, both times.
\end{labbox}

\textbf{Quick mental check.} What happens if you multiply an $m\times n$ matrix $\mat{A}$ by the first standard coordinate vector $\vec{e}_1 = \colvec{1\\0\\ \vdots \\ 0}$? By the column rule,
\[
  \mat{A}\vec{e}_1 = 1\vec{a}_1 + 0\vec{a}_2 + \cdots + 0\vec{a}_n = \vec{a}_1.
\]
Multiplying by $\vec{e}_1$ plucks out the first column of $\mat{A}$. In general, $\mat{A}\vec{e}_j$ extracts the $j$-th column of $\mat{A}$, where $\vec{e}_j$ has a $1$ in position $j$ and $0$ everywhere else. We will use this fact often.

\subsection{What the column picture says about solving $\mat{A}\vec{x} = \vec{b}$}

The equation $\mat{A}\vec{x} = \vec{b}$ asks: \emph{find weights $\vec{x}$ so that the linear combination of the columns equals $\vec{b}$}. In our column language:
\[
  \mat{A}\vec{x} = \vec{b}
  \quad\Longleftrightarrow\quad
  x_1\vec{a}_1 + \cdots + x_n\vec{a}_n = \vec{b}
  \quad\Longleftrightarrow\quad
  \vec{b} \in \spn\{\vec{a}_1,\dots,\vec{a}_n\}.
\]
The double arrow $\Longleftrightarrow$ means ``each statement implies the other'', so the three statements are the same condition written three ways. \textbf{The system $\mat{A}\vec{x}=\vec{b}$ has a solution if and only if $\vec{b}$ lies in the span of the columns of $\mat{A}$.} (``If and only if'' also means that both directions hold: a solution puts $\vec{b}$ in the span, and $\vec{b}$ in the span gives a solution.) The span of the columns is so important it gets a name -- the \textbf{column space} of $\mat{A}$, written $\col(\mat{A})$ -- and an entire chapter later in the book. You are meeting it here as ``the set of all reachable right-hand sides.''

\begin{focusbox}
Solvability is a geometry question. ``Can I solve $\mat{A}\vec{x}=\vec{b}$?'' is \emph{exactly} ``is the target $\vec{b}$ sitting inside the span of the columns?'' If $\vec{b}$ is in the column space, a solution exists. If $\vec{b}$ floats outside it, no choice of weights will ever reach it, and the system has no solution. No cleverness with the algebra can change that; the target is out of reach.
\end{focusbox}

\section{Reaching a Target $\vec{b}$}

A few small cases turn solvability into pictures, and the geometry tells you the answer before any algebra does.

\textbf{Case 1: columns span everything.} Take
\[
  \mat{A} = \begin{bmatrix} 1 & 3 \\ 2 & 1 \end{bmatrix}.
\]
Its columns $\colvec{1\\2}$ and $\colvec{3\\1}$ point in different directions, so they span all of $\R^2$. Therefore \emph{every} target $\vec{b}\in\R^2$ is reachable -- $\mat{A}\vec{x}=\vec{b}$ is solvable for every $\vec{b}$, as we saw in Example~\ref{ch03:ex-in-span}. The column space fills the whole plane; nothing is out of reach.

\textbf{Case 2: columns collapse to a line.} Take
\[
  \mat{B} = \begin{bmatrix} 1 & 2 \\ 2 & 4 \end{bmatrix}.
\]
Its second column $\colvec{2\\4}$ is exactly twice the first column $\colvec{1\\2}$. They lie on the same line. So no matter what weights $x_1, x_2$ you choose,
\[
  x_1\colvec{1\\2} + x_2\colvec{2\\4} = (x_1 + 2x_2)\colvec{1\\2}
\]
is always a multiple of $\colvec{1\\2}$. The column space is a single line -- the line through the origin in direction $\colvec{1\\2}$. The system $\mat{B}\vec{x}=\vec{b}$ is solvable \emph{only} for targets $\vec{b}$ sitting on that line, and unsolvable for every $\vec{b}$ off it, as Figure~\ref{ch03:collapse-line} shows.

\begin{figure}[H]
\centering
\begin{tikzpicture}[scale=0.8,>=Stealth]
  \draw[help lines,gray!22] (-1.4,-1.4) grid (4.4,4.9);
  \draw[->] (-1.4,0)--(4.4,0) node[right]{$x$};
  \draw[->] (0,-1.4)--(0,4.9) node[above]{$y$};
  % column space line: direction (1,2)
  \draw[sourceblue!55,very thick] (-0.7,-1.4)--(2.45,4.9);
  \draw[->,vecblue,very thick] (0,0)--(1,2) node[above left]{$\vec{a}_1$};
  \draw[->,vecred,very thick] (0,0)--(2,4) node[below right=1pt]{$\vec{a}_2=2\vec{a}_1$};
  \node[sourceblue!70!black, anchor=west] at (2.5,4.7) {\small $\col(\mat{B})$};
  % reachable target
  \fill[targetgreen] (1.5,3) circle (2.5pt);
  \node[below right=4pt,targetgreen!60!black] at (1.5,3) {\small reachable $\vec{b}$};
  % unreachable target
  \fill[trapred] (3,1.2) circle (2.5pt);
  \node[right=4pt,trapred] at (3,1.2) {\small unreachable $\vec{b}$};
\end{tikzpicture}
\caption{When the columns are parallel, the column space collapses to a line. Targets on the line (green) are reachable; targets off the line (red) are forever out of reach, no matter the weights.}
\label{ch03:collapse-line}
\end{figure}

\begin{example}[Reading solvability off the geometry]\label{ch03:ex-solvable}
For $\mat{B}=\begin{bmatrix}1&2\\2&4\end{bmatrix}$ above, is $\mat{B}\vec{x}=\colvec{3\\6}$ solvable? Is $\mat{B}\vec{x}=\colvec{3\\1}$ solvable?

\emph{First target.} $\colvec{3\\6} = 3\colvec{1\\2}$ sits right on the line, so yes -- pick $x_1=3, x_2=0$ (or $x_1=1,x_2=1$, or infinitely many others, since $x_1+2x_2=3$ is one equation in two unknowns). The target is reachable in infinitely many ways.

\emph{Second target.} $\colvec{3\\1}$ is \emph{not} a multiple of $\colvec{1\\2}$ (that would need the second entry to be $6$, not $1$). It lies off the line, so no weights reach it and the system has no solution. The algebra confirms it: the two equations $x_1+2x_2=3$ and $2x_1+4x_2=1$ say $2(x_1+2x_2)$ equals both $6$ and $1$, a contradiction.
\end{example}

\begin{trapbox}
A \emph{square} matrix (as many rows as columns) does \emph{not} guarantee a solution. Both $\mat{A}$ and $\mat{B}$ above are $2\times 2$, but $\mat{A}$ solves for every $\vec{b}$ while $\mat{B}$ fails for most. What decides it is whether the columns span the whole target space or collapse onto a smaller flat piece through the origin, such as a line inside the plane. Always ask ``what do the columns span?'' before assuming a solution exists.
\end{trapbox}

\begin{applicationbox}
\textbf{Color, recipes, and recommender systems.} Linear combinations are how machines mix and reach things.

The color of a screen pixel is a linear combination of three fixed color vectors, the red, green and blue ``primaries.'' The weights $(r,g,b)$ are the brightness dials. A pixel cannot shine with negative brightness or above its maximum, so each weight is limited to a range (say $0$ to $1$). The colors a monitor can show, called its \emph{gamut}, are the combinations whose weights stay in that range: a bounded region inside the span, not the whole span. A color outside the gamut cannot be shown on that screen; it is an unreachable target, like the red dot in Figure~\ref{ch03:collapse-line}. Wide-gamut displays use more saturated primaries, which point further apart and enlarge the region.

The same idea appears in recommender systems (the programs that suggest films or songs). One common model, \emph{matrix factorization}, writes the large table of user ratings as the product of two thin tables. In that model each user's predicted ratings, one number per film, form a vector that is a linear combination of a handful of hidden ``taste'' vectors, and the weights describe that user. Every prediction the model can make lies in the span of those taste vectors. The dimension of that span (how many different taste directions the model keeps) is a design choice called the \emph{rank}, and you will meet it properly with the singular value decomposition (SVD) and principal component analysis (PCA) later in the book.
\end{applicationbox}

\section{When a New Vector Adds Nothing: Dependence, Previewed}

Sometimes a new vector enlarges the span, and sometimes it adds nothing. The full theory comes later; here we name the pattern.

\begin{definition}[Adds nothing / redundant -- informal]\label{ch03:def-redundant}
A vector $\vec{v}$ is \textbf{redundant} for a collection if it already lies in the span of the others -- that is, $\vec{v}$ is itself a linear combination of them. Adding a redundant vector does not change the span.
\end{definition}

Two quick litmus tests you can run by eye:

\begin{itemize}
  \item \textbf{One vector is a multiple of another} ($\vec{v}_2 = c\,\vec{v}_1$): redundant. They share a line. This is the case of $\mat{B}$ above.
  \item \textbf{One vector lies in the plane of two others} ($\vec{v}_3 = a\vec{v}_1 + b\vec{v}_2$): redundant. It adds no new direction out of that plane.
\end{itemize}

When \emph{no} vector in the set is redundant -- when every one of them points somewhere the others cannot reach -- we call the set \textbf{linearly independent}. When at least one is redundant, the set is \textbf{linearly dependent}. Figure~\ref{ch03:redundancy-r2} shows a dependent set of three vectors in the plane.

\begin{figure}[H]
\centering
\begin{tikzpicture}[scale=0.9,>=Stealth]
  \draw[help lines,gray!25] (-1.9,-0.9) grid (3.9,3.9);
  \draw[->] (-1.9,0)--(3.9,0) node[right]{$x$};
  \draw[->] (0,-0.9)--(0,3.9) node[above]{$y$};
  % v1 and v2
  \draw[->,vecblue,very thick] (0,0)--(2,1) node[below right]{$\vec{v}_1$};
  \draw[->,vecred,very thick]  (0,0)--(-1,2) node[above left]{$\vec{v}_2$};
  % parallelogram decomposition of v3 = v1 + v2 = (1,3)
  \draw[vecblue!40,thick,dashed] (-1,2)--(1,3);
  \draw[vecred!40,thick,dashed]  (2,1)--(1,3);
  % v3 arrow
  \draw[->,structurepurple,very thick] (0,0)--(1,3) node[above right]{$\vec{v}_3 = \vec{v}_1 + \vec{v}_2$};
  \fill[structurepurple] (1,3) circle (2.5pt);
\end{tikzpicture}
\caption{Redundancy in $\R^2$. Two non-parallel vectors $\vec{v}_1$ and $\vec{v}_2$ already span the entire plane. Any third vector $\vec{v}_3$ lies in that plane and can be built as a linear combination of them ($1\vec{v}_1 + 1\vec{v}_2$ here). Because $\vec{v}_3$ points nowhere new, it is redundant -- the set $\{\vec{v}_1,\vec{v}_2,\vec{v}_3\}$ is linearly dependent.}
\label{ch03:redundancy-r2}
\end{figure}

\begin{theorem}[Two non-parallel vectors span all of $\R^2$]\label{ch03:thm-span-r2}
Let $\vec{v}_1,\vec{v}_2 \in \R^2$ be nonzero and non-parallel (neither is a scalar multiple of the other). Then $\spn\{\vec{v}_1,\vec{v}_2\} = \R^2$: every vector in the plane is some linear combination of them.
\end{theorem}

\begin{proof}
Write $\vec{v}_1 = \colvec{a\\c}$ and $\vec{v}_2 = \colvec{b\\d}$, and let $\vec{b}^\star = \colvec{p\\q}$ be \emph{any} target (the star only keeps the target's name apart from the number $b$). We must find weights $x_1,x_2$ with $x_1\vec{v}_1 + x_2\vec{v}_2 = \vec{b}^\star$, i.e.
\[
  \begin{aligned}
    a x_1 + b x_2 &= p,\\
    c x_1 + d x_2 &= q.
  \end{aligned}
\]
Solve by elimination. Multiply the first equation by $d$, the second by $b$, and subtract:
\[
  (ad - bc)\,x_1 = dp - bq.
\]
The number $ad-bc$ is the \emph{determinant} of the $2\times 2$ matrix with columns $\vec{v}_1,\vec{v}_2$. Its size is the area of the parallelogram with sides $\vec{v}_1$ and $\vec{v}_2$, and its sign ($+$ or $-$) records whether $\vec{v}_2$ is reached from $\vec{v}_1$ by turning counterclockwise or clockwise (Figure~\ref{ch03:det-area}).

We need $ad-bc\neq 0$, and the non-parallel hypothesis gives it. Suppose instead that $ad-bc=0$, that is, $ad=bc$. If $a\neq 0$, put $\lambda=b/a$; then $b=\lambda a$ and $d=bc/a=\lambda c$, so $\vec{v}_2=\lambda\vec{v}_1$. If $a=0$, then $bc=0$; since $\vec{v}_1\neq\vec{0}$ we have $c\neq 0$, so $b=0$ and $\vec{v}_2=\colvec{0\\d}=\tfrac{d}{c}\vec{v}_1$. Either way the vectors are parallel, which the hypothesis rules out. So $ad-bc\neq 0$, and we may divide:
\[
  x_1 = \frac{dp - bq}{ad-bc}.
\]
In the same way, multiplying the first equation by $c$, the second by $a$, and subtracting gives $(ad-bc)\,x_2 = aq-cp$, so
\[
  x_2 = \frac{aq - cp}{ad-bc}.
\]
Substituting these two values back into both equations confirms that they hold. These weights exist for \emph{every} $(p,q)$, so every target is reachable. Hence the span is all of $\R^2$.
\end{proof}

\begin{figure}[htbp]
\centering
\begin{tikzpicture}[scale=0.85,>=Stealth]
  % left: non-parallel
  \begin{scope}
    \draw[help lines,gray!25] (-0.4,-0.4) grid (3.9,3.4);
    \draw[->] (-0.4,0)--(3.9,0) node[right]{$x$};
    \draw[->] (0,-0.4)--(0,3.4) node[above]{$y$};
    \fill[targetgreen!20] (0,0)--(2,1)--(3,3)--(1,2)--cycle;
    \draw[->,vecblue,very thick] (0,0)--(2,1) node[right]{$\vec{v}_1$};
    \draw[->,vecred,very thick] (0,0)--(1,2) node[left]{$\vec{v}_2$};
    \node[targetgreen!50!black] at (1.55,1.5) {\small area $3$};
    \node[align=center] at (1.75,-1.3) {\small $\vec{v}_1=(2,1)$, $\vec{v}_2=(1,2)$\\ \small $ad-bc=2\cdot 2-1\cdot 1=3\neq 0$};
  \end{scope}
  % right: parallel
  \begin{scope}[xshift=6.6cm]
    \draw[help lines,gray!25] (-0.4,-0.4) grid (3.9,3.4);
    \draw[->] (-0.4,0)--(3.9,0) node[right]{$x$};
    \draw[->] (0,-0.4)--(0,3.4) node[above]{$y$};
    \draw[->,vecred,very thick] (0,0)--(3,1.5) node[above]{$\vec{v}_2$};
    \draw[->,vecblue,very thick] (0,0)--(2,1) node[below right]{$\vec{v}_1$};
    \node[align=center] at (1.75,-1.3) {\small $\vec{v}_1=(2,1)$, $\vec{v}_2=(3,1.5)$\\ \small $ad-bc=2\cdot 1.5-3\cdot 1=0$};
  \end{scope}
\end{tikzpicture}
\caption{The determinant $ad-bc$ as an area. Left: two non-parallel vectors enclose a parallelogram of area $3$, so $ad-bc\neq 0$. Right: when $\vec{v}_2$ is a multiple of $\vec{v}_1$ the parallelogram is squashed flat, its area is $0$, and the two vectors reach only one line.}
\label{ch03:det-area}
\end{figure}

\begin{remark}
The proof links four statements: ``the vectors are non-parallel'' (geometry), ``$ad-bc\neq 0$'' (algebra), ``the columns span $\R^2$'' (span), and ``$\mat{A}\vec{x}=\vec{b}$ is solvable for every $\vec{b}$'' (linear systems) all state the exact same condition. Much of linear algebra consists of translating fluently among these four views.
\end{remark}

\subsection{Counting: why you cannot span $\R^3$ with two vectors}

The \emph{dimension} of a span is its number of independent directions: $0$ for the origin alone, $1$ for a line, $2$ for a plane, $3$ for all of $\R^3$ (the rungs of Figure~\ref{ch03:ladder}).

\begin{proposition}[Too few vectors cannot span]\label{ch03:prop-too-few}
The span of $k$ vectors in $\R^n$ is at most $k$-dimensional. In particular, fewer than $n$ vectors can never span all of $\R^n$.
\end{proposition}

We will only sketch the reasoning now (the airtight proof waits for the dimension chapter). Each vector you add can push the span up by \emph{at most} one dimension -- one new direction at most. Starting from nothing (dimension $0$), after $k$ vectors you have climbed at most $k$ rungs of the point-line-plane-space ladder. So two vectors span at most a plane; they sit flat inside $\R^3$ and can never fill it. You need at least three vectors to span $\R^3$, at least $n$ to span $\R^n$. Conversely, more vectors than necessary guarantees redundancy: any $n+1$ vectors in $\R^n$ \emph{must} be dependent, because the $(n{+}1)$-th has nowhere new to point.

\textbf{Quick mental check.} Can four vectors in $\R^3$ ever be linearly independent? No. If the first three are independent, they already fill all of $3$-space, so the fourth has nowhere new to point and is redundant. If the first three are not independent, the set is dependent already.

\begin{connectionbox}
\textbf{Spans show up far past arrows in $\R^n$.} Once you recognize the pattern ``set of all combinations of a few generators,'' you start seeing it everywhere.
\begin{itemize}
  \item \textbf{Polynomials.} The functions $1, t, t^2$ are ``vectors,'' and their span is every quadratic $a + bt + ct^2$. Calculus students compute in this span constantly without naming it. A Taylor polynomial of degree $n$ (the polynomial approximation to a function that calculus builds from its derivatives at one point) lies in the span of $1, t, \dots, t^n$.
  \item \textbf{Differential equations.} The solutions of $y'' + y = 0$ (where $y''$ is the second derivative of the unknown function $y(t)$) are exactly the linear combinations $c_1\cos t + c_2\sin t$: a $2$-dimensional span of two building-block solutions. Solving a linear differential equation of this kind means finding such building blocks and taking their span.
  \item \textbf{Signals and Fourier.} A sound wave is (approximately) a linear combination of pure sine tones, an idea due to Joseph Fourier. The span of the basic tones is the space of signals you can build; an audio equalizer lets you turn the weight dials.
  \item \textbf{Geometry of subspaces.} A line through the origin, a plane through the origin, and all of $\R^3$ are all examples of the same construction: spans of $1$, $2$, and $3$ independent vectors. Taking spans produces \emph{subspaces}, the subject of the middle chapters of this book.
\end{itemize}
Whenever you see ``all combinations of a few building blocks,'' you are looking at a span, and everything in this chapter applies.
\end{connectionbox}

\begin{strangbox}
\textbf{Before Strang:} when he opens his course with the ``column picture'' and insists $\mat{A}\vec{x}$ is a combination of columns, you will already own it -- you derived it above by regrouping $\colvec{c_1+3c_2\\2c_1+c_2}$ into columns. His famous line, ``linear algebra is about $\mat{A}\vec{x}=\vec{b}$, and the question is whether $\vec{b}$ is a combination of the columns,'' is precisely the chain we built in the section on reaching a target.

\textbf{After Strang:} when you return for the \emph{column space} $\col(\mat{A})$ and the \emph{four fundamental subspaces}, recognize that the column space is the span of the columns introduced here, studied as a subspace in its own right. The rank of a matrix is the dimension of this span: the number of independent directions among its columns. Every later theorem about rank is a theorem about the dimension of a span.
\end{strangbox}

\begin{aestheticbox}
Two operations, scaling and adding, are enough to produce the lines, planes and higher-dimensional
flat pieces of the subject, and the test for whether $\mat{A}\vec{x}=\vec{b}$ has a solution.
Linear independence, dimension and rank are all built on them.
\end{aestheticbox}

\begin{historybox}
The idea that spaces of any dimension can be generated by taking linear combinations of a few
directions belongs above all to Hermann Grassmann (1809--1877), a secondary school teacher
in Stettin, Prussia. In the 1830s, mathematicians and physicists were stuck in a computational bottleneck:
problems in tidal forces, rotation, and mechanics were buried under mountains of coordinate calculations.
If you rotate your coordinate axes or shift your origin, every single $(x,y,z)$ number changes, even though
the physical object and the forces acting on it have not changed at all. The only available tools were
Cartesian coordinates (grids of numbers pegged to fixed perpendicular axes, introduced by Ren\'e Descartes)
and classical Euclidean geometry, which had no language for spaces beyond three dimensions. Grassmann wanted
an algebra of geometry that worked directly on directed geometric quantities in four, five, or $n$ dimensions
without coordinates.

In his book of 1844, \emph{Die Lineare Ausdehnungslehre} (``The Theory of Linear Extension''), Grassmann
showed that linear combinations $c_1 \vec{v}_1 + \cdots + c_k \vec{v}_k$ sweep out flat linear spaces:
lines, planes, and higher-dimensional hyperplanes (flat slices of space that have one dimension less than the
ambient space). Because scaling and adding never bend anything, these sets are always flat.

Grassmann's book was decades ahead of its time, written in dense, philosophical German without numerical
worked examples. He hoped it would win him a university post; it did not, and its readers found it
unreadable. Even August Ferdinand M\"obius (the Leipzig mathematician best known
for the one-sided M\"obius strip) confessed he could not follow Grassmann's terminology. The mathematical
community refused to believe that algebraic operations could be performed on geometric magnitudes without
coordinates. Disheartened by decades of obscurity, Grassmann turned in mid-life from mathematics to
historical linguistics, where he became famous for a large dictionary of the Rigveda and for
Grassmann's Law, a rule about sound changes in Sanskrit and ancient Greek. Only in 1888, eleven
years after Grassmann's death, did the Italian mathematician Giuseppe Peano extract the modern axioms of
vector spaces and linear combinations from Grassmann's work.

Grassmann's notion of span gave linear systems a geometric formulation: solving $\mat{A}\vec{x}=\vec{b}$
means asking whether the target $\vec{b}$ lies inside the subspace spanned by the columns of $\mat{A}$.
Today the research questions concern scale. When a dataset has millions of
columns, testing whether a combination exists is routine, but two harder problems are areas of active research.
The first is \emph{sparse recovery}: finding the recipe that hits or approximates a target $\vec{b}$ using the fewest
nonzero weights. It is the central problem of \emph{compressed sensing}, which rebuilds signals such as MRI
scans from fewer measurements than usual, and in general it is \emph{NP-hard}, meaning no method is known that
solves every case quickly as the problem grows.
The second is \emph{lattice problems}: finding the shortest nonzero vector when weights are restricted to integers rather than
real numbers. These are believed to be so hard to solve that \emph{post-quantum cryptography}, encryption designed to
resist future quantum computers, rests on them.
\end{historybox}

\section{Summary}

\begin{itemize}
  \item A \textbf{linear combination} $c_1\vec{v}_1 + \cdots + c_k\vec{v}_k$ is built from only two moves: \emph{scaling} (along a direction) and \emph{adding} (mixing directions). The $c_i$ are weights.
  \item The \textbf{span} is the set of \emph{all} linear combinations -- every destination the weights can reach. It always passes through the origin.
  \item Geometrically the span climbs a ladder: \emph{point} (zero vectors), \emph{line} (one direction), \emph{plane} (two independent directions), \emph{space}, and on up. Each new vector adds at most one rung.
  \item A vector that already lies in the others' span is \textbf{redundant}; it adds nothing. ``No redundancy'' is \textbf{linear independence}; this is the geometric seed of independence, rank, and dimension.
  \item \textbf{Column picture:} $\mat{A}\vec{x} = x_1\vec{a}_1 + \cdots + x_n\vec{a}_n$ is a linear combination of the columns of $\mat{A}$ with weights $\vec{x}$.
  \item Therefore $\mat{A}\vec{x}=\vec{b}$ is solvable \emph{exactly when} $\vec{b}$ lies in the span of the columns (the \textbf{column space}). Solvability is a geometry question.
\end{itemize}

\begin{recoverybox}
\textbf{Cost and intuition recap.} The \emph{intuition}: a span is everywhere you can get to by scaling and adding, and it grows by whole rungs -- point, line, plane, space. The \emph{cost}: to test whether a target $\vec{b}$ lies in $\spn\{\vec{v}_1,\dots,\vec{v}_k\}$, you do not search through weights; you set up one linear system, $c_1\vec{v}_1+\cdots+c_k\vec{v}_k=\vec{b}$ ($n$ equations, $k$ unknowns), and solve it once. ``Is $\vec{b}$ in the span?'' costs the same as ``does this system have a solution?''
\end{recoverybox}

\subsection*{Looking ahead: the missing ruler and protractor}

We now have a complete geometric picture of reachability: scaling and adding sweeps out flat spans, and solving $\mat{A}\vec{x}=\vec{b}$ asks whether the target $\vec{b}$ lands on that flat sheet.

Something is missing. Suppose $\vec{b}$ does \emph{not} lie in the span -- like the unreachable red dot in Figure~\ref{ch03:collapse-line}, or an \emph{overdetermined} system in data science (more equations than unknowns) where measurements are corrupted by noise. What do we do then?

In practice, a target is almost never perfectly reachable. If we cannot reach $\vec{b}$ exactly, the natural engineer's question is: \emph{which reachable point comes closest?} And how far away is it?

\begin{figure}[H]
\centering
\begin{tikzpicture}[scale=1.3,>=Stealth,
   x={(0.95cm,0.25cm)}, y={(-0.8cm,0.45cm)}, z={(0cm,1.1cm)}]
  \fill[targetgreen!12] (-1.3,-1.2,0)--(1.7,-1.2,0)--(1.7,1.3,0)--(-1.3,1.3,0)--cycle;
  \draw[targetgreen!50,thick] (-1.3,-1.2,0)--(1.7,-1.2,0)--(1.7,1.3,0)--(-1.3,1.3,0)--cycle;
  \node[targetgreen!70!black] at (0.45,-0.7,0) {\small $\spn\{\vec{a}_1,\vec{a}_2\} = \col(\mat{A})$};
  \draw[->,vecblue,very thick] (0,0,0)--(1.2,0.2,0) node[right]{$\vec{a}_1$};
  \draw[->,vecblue,very thick] (0,0,0)--(0.2,1.0,0) node[left]{$\vec{a}_2$};
  \coordinate (P) at (1.1,0.7,0);
  \draw[->,targetgreen,very thick] (0,0,0)--(P);
  \fill[targetgreen] (P) circle (2pt) node[right=3pt,targetgreen]{$\vec{p}$};
  \coordinate (B) at (1.1,0.7,1.4);
  \draw[->,vecred,very thick] (0,0,0)--(B) node[above]{$\vec{b}$};
  \fill[vecred] (B) circle (2pt);
  \draw[->,trapred,thick,dashed] (P)--(B) node[midway,right]{$\vec{e} = \vec{b} - \vec{p}$};
  \draw[gray!70,thin] ($(P)+(0.15,0,0)$)--($(P)+(0.15,0,0.15)$)--($(P)+(0,0,0.15)$);
\end{tikzpicture}
\caption{The geometric puzzle that motivates Chapter~4 (Section~\ref{ch04:dot-def}). When a target $\vec{b}$ (red) floats outside the column space $\col(\mat{A})$ (green plane), no linear combination can reach it: $\mat{A}\vec{x}=\vec{b}$ has no exact solution. The best reachable substitute is the projection $\vec{p}$ (green), which drops a perpendicular shadow onto the plane so that the error $\vec{e} = \vec{b}-\vec{p}$ is as short as possible. But measuring lengths and perpendicular angles requires an operation that addition and scaling cannot provide -- the dot product.}
\label{ch03:projection-bridge}
\end{figure}

Addition and scaling alone cannot answer that question. They build the flat stage, but they come with neither a ruler to measure distances nor a protractor to measure angles. To find the closest point $\vec{p}$ on the span to $\vec{b}$, we must drop a perpendicular line from $\vec{b}$ down onto the plane (Figure~\ref{ch03:projection-bridge}). But ``perpendicular'' is an angle concept, and ``closest'' is a length concept. Neither exists until we introduce a way to measure them.

In Chapter~4 (Section~\ref{ch04:dot-def}) we introduce the single operation that supplies both the ruler and the protractor: the \textbf{dot product}: multiply matching entries and add them up. From that small sum come lengths, angles, perpendicularity, and the projections that give the best approximate solution of a system that has no exact solution.

\section{Exercises}

\begin{exercise}[Warmup combinations]\label{ch03:ex-warmup}
Let $\vec{u} = \colvec{1\\-1}$ and $\vec{w} = \colvec{2\\3}$. Compute each linear combination:
(a) $3\vec{u} + \vec{w}$;\;
(b) $2\vec{u} - 2\vec{w}$;\;
(c) $0\vec{u} + 5\vec{w}$;\;
(d) $-\vec{u} - \vec{w}$.
For (c), what does the weight $0$ on $\vec{u}$ tell you about where the result lands?
\end{exercise}

\begin{exercise}[Solve for the weights]\label{ch03:ex-find-weights}
Find weights $c_1, c_2$ so that $c_1\colvec{1\\1} + c_2\colvec{1\\-1} = \colvec{4\\2}$. Then verify your answer by substituting back. (Hint: add and subtract the two scalar equations.)
\end{exercise}

\begin{exercise}[In the span or not?]\label{ch03:ex-membership}
Let $\vec{v}_1 = \colvec{2\\1}$ and $\vec{v}_2 = \colvec{-4\\-2}$.
(a) What is $\spn\{\vec{v}_1,\vec{v}_2\}$ geometrically -- a point, a line, or a plane? Why?
(b) Is $\colvec{6\\3}$ in the span? Is $\colvec{1\\0}$ in the span? Justify each answer geometrically, then confirm algebraically.
\end{exercise}

\begin{exercise}[Column picture by hand]\label{ch03:ex-column}
Let $\mat{A} = \begin{bmatrix} 1 & 0 & 2 \\ 0 & 1 & 3 \end{bmatrix}$ and $\vec{x} = \colvec{4\\-1\\2}$. Compute $\mat{A}\vec{x}$ as a linear combination of the columns of $\mat{A}$ (write out ``$4$ times column one, plus \dots''). Then state what $\spn$ of the three columns is, and explain why the third column is redundant.
\end{exercise}

\begin{exercise}[Reaching a target]\label{ch03:ex-target}
For $\mat{A} = \begin{bmatrix} 1 & 2 \\ 3 & 6 \end{bmatrix}$:
(a) Describe $\col(\mat{A})$ (the span of the columns) as a specific line through the origin.
(b) Is $\mat{A}\vec{x} = \colvec{2\\6}$ solvable? Is $\mat{A}\vec{x} = \colvec{2\\5}$ solvable? Answer by checking whether each target lies on the line from (a).
(c) For the solvable case, find \emph{all} solutions $\vec{x}$, and explain geometrically why there are infinitely many.
\end{exercise}

\begin{exercise}[The third vector decides]\label{ch03:ex-third}
In $\R^3$, let $\vec{v}_1 = \colvec{1\\0\\0}$ and $\vec{v}_2 = \colvec{0\\1\\0}$.
(a) Describe $\spn\{\vec{v}_1,\vec{v}_2\}$ as a specific plane.
(b) For $\vec{v}_3 = \colvec{2\\3\\0}$, what is $\spn\{\vec{v}_1,\vec{v}_2,\vec{v}_3\}$? Did $\vec{v}_3$ add a new direction?
(c) For $\vec{v}_3 = \colvec{2\\3\\1}$, what is $\spn\{\vec{v}_1,\vec{v}_2,\vec{v}_3\}$ now? What changed?
\end{exercise}

\begin{exercise}[A span that misses the origin?]\label{ch03:ex-affine}
Your friend claims the line $y = x + 1$ in the plane is the span of some single vector. Prove your friend is wrong in one sentence. (Hint: which point is in \emph{every} span, and is it on this line?)
\end{exercise}

\begin{exercise}[Counting argument]\label{ch03:ex-counting}
Explain, in your own words and with a picture, why three vectors in $\R^2$ \emph{must} be linearly dependent (at least one is redundant). Then exhibit explicit weights, not all zero, with $c_1\colvec{1\\0} + c_2\colvec{0\\1} + c_3\colvec{1\\1} = \colvec{0\\0}$, demonstrating the dependence concretely.
\end{exercise}

\begin{exercise}[Same span, different generators]\label{ch03:ex-same-span}
Show that $\spn\left\{\colvec{1\\0},\colvec{0\\1}\right\} = \spn\left\{\colvec{1\\1},\colvec{1\\-1}\right\}$ by proving each side contains the other. (The symbol $\subseteq$ means ``is contained in.'' To show one span is contained in the other, write a general vector of the first span as a combination of the other set's vectors.) What does this say about the idea that a span has many possible ``generating sets''?
\end{exercise}

\begin{exercise}[Proof: span is closed under combination]\label{ch03:ex-closed}
Let $S = \spn\{\vec{v}_1,\dots,\vec{v}_k\}$. Suppose $\vec{p}$ and $\vec{q}$ are both in $S$. Prove that $\alpha\vec{p} + \beta\vec{q}$ is also in $S$ for any scalars $\alpha,\beta$. (Write $\vec{p}$ and $\vec{q}$ as combinations of the $\vec{v}_i$, then combine.) This ``closed under linear combination'' property is the defining feature of a \emph{subspace} -- you have just proved every span is one.
\end{exercise}

\begin{exercise}[Challenge: build any vector two ways]\label{ch03:ex-challenge}
Let $\vec{v}_1 = \colvec{1\\0}$, $\vec{v}_2 = \colvec{0\\1}$, $\vec{v}_3 = \colvec{1\\1}$.
(a) Find two \emph{different} sets of weights $(c_1,c_2,c_3)$ that both produce $\colvec{3\\2}$.
(b) Explain why having a redundant vector in the set is exactly what makes more than one set of weights possible. Connect this to why $\mat{B}\vec{x}=\colvec{3\\6}$ in Example~\ref{ch03:ex-solvable} had infinitely many solutions.
\end{exercise}
